\chapter{Elliptische Kurven über endlichen Körpern}
Für die algebraischen Grundlagen wird \cite{karpf24}, S.317-318 verwendet.
\section{Motivation}
Nicht jede Art von elliptischer Kurve macht Sinn im Bereich der Kryptography. Zum Beispiel würden Kurven über reellen Zahlen Gleitpunkt Arithmetik benötigen, welche Probleme wie Ungenauigkeit oder verringerte Rechenleistung mit sich bringen würden. Daher werden elliptische Kurven über endlichen Körpern definiert.
Hierbei wird wie bereits eingeführt die kurze Weierstraß-Gleichung verwendet(vgl. \cite{omondi20}, S.230). Man benötigt für Kapitel 5 sichere Kurven, daher muss man gewisse Dinge über die Ordnung und Eigenschaften erfahren.

\section{Endliche Körper}
\begin{definition}
    Ein endlicher Körper ist ein Körper mit endlich vielen Elementen. Für jede Primzahlpotenz \(q=p^n\) existiert bis auf Isomorphie genau ein endlicher Körper mit \(q\) Elementen. Dieser wird dann mit \(\mathbb F_q\) oder \(GF(q)\)(Galois Field) bezeichnet.
\end{definition}
In kryptographischen Anwendungen treten insbesondere die Körper \(\mathbb F_p\) und \(\mathbb F_{2^m}\) auf. Der kleinste Durchscnitt $P$ aller Teilkörper von $K$ wird der \textbf{Primkörper} genannt
\[
    P=\cap \{U|U \text{ist ein Teilkörper von }K\}.
\]
\begin{satz}

    \label{hilfs:primkörper}
    Seien $K$ ein Körper und $K_0$ der Primkörper von $K$.
    \begin{enumerate}
        \item Ist char($K)=0$, so ist $K_0$ zu $\mathbb{Q}$ isomorph.
        \item Ist char($K)=p\neq 0$, so ist $K_0$ zu $\mathbb{Z}/p\cdot \mathbb{Z}$ isomorph.
    \end{enumerate}
\end{satz}
Mit Worten heißt dies: Hat $K$ die Charakteristik 0, so ist der Primkörper von $K$ zu $\mathbb{Q}$ isomorph. Ist die Charakteristik $p \neq 0$, so ist er zu $\mathbb{Z}/p$ isomorph.
\begin{proof}
    \begin{enumerate}
        \item Sei \(\Phi:\mathbb{N}\leftarrow K_0\) eine Abbildung für alle $n\in\mathbb{N}$. Sei
              \[
                  n^\Phi = (1_k+\cdots+1_k). \quad \text{n-mal}.
              \]
              Sind $n,m\in\mathbb{N}$ und $n\leq m$ mit der Eigenschaft $n^\psi=m^\psi$, so folgt
              \[
                  0_K=n^\psi-m^\psi=(1_K+\cdots+1_K). \quad \text{(n-m)-mal}.
              \]
              Da char($K)=0$. ist ist, dann $n=m$. Also ist $\psi$ eine injektive Abbildung. Somit besitzt jede $K_0$ eine Teilmenge, die sich bezüglich der Addition wie $\mathbb{N}$ verhält. Da $K_0$ ein Körper ist, gibt es also in $K_0$ einen zu $\mathbb{Z}$ isomorphen Teilring und dann auch einen zu $\mathbb{Q}$ isomorphen Teilkörper. Da $\mathbb{Q}$ keine echten Teilkörper hat, ist $K_0$ selbst isomorph zu $\mathbb{Q}$.
        \item Seien nun $p$ eine Primzahl und char($K)= p$. Sei weiter
              \[
                  L = {0_K,1_K,\dots,1_k+\cdots+1_K}\quad \text{(p-1)-mal}
              \]
              Da char($K)=p$ ist, sind die Elemente in $L$ paarweise verschieden und daher ist $|L|=p$, außerdem ist $L$ unter Addition abgeschlossen.

    \end{enumerate}
\end{proof}
\cite{stroth24}, S.66--67
\begin{satz}
    Sei $K$ ein endlicher Körper. Dann gibt es eine Primzahl $p$ und ein $n\in \mathbb{N}$, sodass \(|K|=p^n\) ist.
\end{satz}

\begin{proof}
    Da \(K\) endlich ist, müssen in der Folge
    \[
        0,\ 1_K,\ 1_K+1_K,\ 1_K+1_K+1_K,\ \dots
    \]
    zwei Elemente übereinstimmen. Daher gibt es eine kleinste positive
    natürliche Zahl \(p\) mit
    \[
        p\cdot 1_K=0.
    \]
    Diese Zahl ist die Charakteristik von \(K\).

    Die Zahl \(p\) ist eine Primzahl. Wäre nämlich \(p=rs\) mit
    \(1<r,s<p\), so wäre
    \[
        (r\cdot 1_K)(s\cdot 1_K)=p\cdot 1_K=0.
    \]
    Da ein Körper keine Nullteiler besitzt, müsste
    \(r\cdot 1_K=0\) oder \(s\cdot 1_K=0\) gelten. Dies widerspricht
    der Minimalität von \(p\).

    Der von \(1_K\) erzeugte Primkörper besitzt somit \(p\) Elemente
    und ist isomorph zu \(\mathbb{F}_p\). Folglich ist \(K\) ein
    Vektorraum über \(\mathbb{F}_p\). Da \(K\) endlich ist, besitzt
    dieser Vektorraum eine endliche Dimension \(n\).

    Bezüglich einer Basis besitzt jedes Element von \(K\) eine
    eindeutige Darstellung
    \[
        a_1v_1+\dots+a_nv_n
        \qquad\text{mit }a_i\in\mathbb{F}_p.
    \]
    Für jeden der \(n\) Koeffizienten gibt es \(p\) Möglichkeiten.
    Daher gilt
    \[
        |K|=p^n.
    \]
\end{proof}

\cite{stroth24},S.163

\(|K|=p^n\), wenn $n$ der Grad von $K$ über dem nach Lemma 21.1 (b) zu $\mathbb{Z}/p$ isomorphen Primkörper.

\begin{lemma}
    Für einen endlichen Körper K der Charakteristik p gilt und $q=p^n$
    \[
        a^q=a
    \] für jedes $a\in K$.
\end{lemma}

\begin{proof}
    Für \(a=0\) ist die Aussage unmittelbar erfüllt, da $0^{p^n}=0$. Sei daher \(a\in K^\times =K\setminus\{0\}\).
    Wegen $|K^\times|=p^{n-1}$ gilt nach dem kleinen Satz von Fermat vgl. \cite{karpf24} S.46 gilt  in einer endlichen Gruppe $G$, \(a^{|G|} = e\) für alle \(a\in G\). Dann \(a^{p^n-1}=1\) für jedes $a\in K^\times$. Daraus folgt dann $a^{p^n}=a$ für alle $a\in K $.
\end{proof}

(vgl. \cite{karpf24}, S.391--392)
Diese Aussage kann nun auch noch in beide Richtungen formuliert werden, da wir sie später bei einem Beweis benötigen.
\begin{lemma}
    Für \(a\in \overline{\mathbb{F}}_q\) gilt
    \[
        a^q=a\Leftrightarrow a\in\mathbb{F}_q
    \]
\end{lemma}
\begin{proof}
    Nach dem vorherigen Lemma gilt für jedes
    \(a\in\mathbb F_q\), dass \(a^q=a\), also
    \(a^q-a=0\). Somit sind die \(q\) Elemente von
    \(\mathbb F_q\) verschiedene Nullstellen des
    Polynoms \(f(X)=X^q-X\).

    Da \(f\) den Grad \(q\) besitzt, hat es auch über
    dem algebraischen Abschluss \(\overline{\mathbb F}_q\)
    höchstens \(q\) verschiedene Nullstellen. Seine
    Nullstellen sind daher genau die Elemente von
    \(\mathbb F_q\). Folglich gilt für jedes
    \(a\in\overline{\mathbb F}_q\):
    \[
        a^q=a \quad\Longleftrightarrow\quad a\in\mathbb F_q.
    \]
\end{proof}


\begin{definition}[Elliptische Kurve über \(\mathbb F_p\)]
    Sei \(p>3\) eine Primzahl. Eine elliptische Kurve über \(\mathbb F_p\) ist die Menge von allen Punkten $(x,y)\in \mathbb{F}_p$ mit
    \[
        \mathbb F_p = y^2 \equiv x^3+ax+b \pmod p
    \]
    und einem Punkt im Unendlichen.
    \[
        4a^3+27b^2 \not\equiv 0 \pmod p
    \] gilt. Die Bedingung stellt wie im allgemeinen Fall sicher, dass die Kurve nichtsingulär ist. Die Einschränkung $p>3$ erlaubt die Verwendung der kurzen Weierstraß-Gleichung. Für Körper der Charakteristik \(2\) und \(3\) haben wir bereits in Kapitel 3 gesehen, dass die allgemeine Weierstraß-Gleichung verwendet werrden muss.
\end{definition}
(\cite{guneysu24},S.279f.)
\begin{theorem}
    Sei $\mathbb{F}_q$ ein endlicher Körper, dann ist die Gruppe der Punkte $E(\mathbb{F}_q)$ immer eine zyklische Gruppe oder das Produkt von zwei zyklischen Gruppen.
\end{theorem}

\cite{silverman06}, S.29

\begin{proof}

\end{proof}

\section{Die Punktmenge $E(\mathbb{F}_q)$}

\subsection{Explizite Punktbestimmung}
\cite{washington08},S.95, hier noch eigene Beispiele erstellen
\begin{enumerate}
    \item Sei E die Kurve $y^2=x^3+x+1$ über $\mathbb{F}_5$. Um nun die Punkte von E zu erhalten, berchnen wir zuerst $x^3+x+1\mod 5$ und anschließend die quadratwurzel von diesem Polynom über $F_5$.
          \begin{center}
              \begin{tabular}{c|c|c|c}
                  x & $x^3+x+1$ & y       & Points       \\
                  \hline
                  0 & 1         & $\pm$ 1 & (0,1), (0,4) \\
                  1 & 3         & -       & -            \\
                  2 & 1         & $\pm$ 1 & (2,1),(2,4)  \\
                  3 & 1         & $\pm 1$ &              \\
                  4 & 4         & $\pm 2$ & (4,2),(4,3)
              \end{tabular}
          \end{center}
          Beachte: $-1\equiv 4 \mod 5$ und $ -2\equiv 3 \mod 5$\\
          Für $y^2=3$ gibt es keine Wurzel in $F_5$\\

          Also ist nun die Ordnung:
          \[
              \mid E(\mathbb{F}_5)\mid = \# \{(x,y)\in \mathbb{F}_5 \mid y^2=x^3+x+1\}+1 = 9
          \]
          Wir berechnen nun einen Punkt  $(3,1)+(2,4)$ auf $E$:
          \[
              \frac{4-1}{2-3}=\frac{3}{-1}=\frac{3}{4}=3\cdot4^{-1} =3\cdot 4=12  \equiv 2 \mod 5
          \]
          Es ist also $y=2(x-3)+1\equiv 2x$. Wenn wir dass nun in die Gleichung für die Kurve einsetzen erhalten wir
          \[(2x)^2=x^3+x+1\Leftrightarrow 0 = x^3-4x^2+x+1\]. Wir kennen schon zwei Schnittpunkte der Kurve $x=3$ und $x=2$, da $(3,1)$ und $(2,4)$ auf der Gerade und Kurve liegen. Da das Polynom kubisch ist, muss die Summe der drei Nullstellen 4 sein. Damit gilt,
          \[
              3+2+r \equiv 4 \mod 5 \Leftrightarrow r\equiv 4 \mod 5
          \]
          Also ist die dritte Koordinate $x=4$, nun erhalten wir $y=2x\equiv 3$, gespiegelt über die x-Achse erhält man die Summe
          \[(3,1)+(2,4)=(4,2)\]

\end{enumerate}




\begin{theorem}
    Sei $E$ eine elliptische Kurve über einem endlichen Körper $\mathbb{F}_q$. Dann gilt,
    \[
        E(\mathbb{F}_q) \simeq \mathbb{Z}_n \text {  oder  } \mathbb{Z}_{n_1} \oplus \mathbb{Z}_{n_2}
    \]
    für ein $n \geq 1$ und für beliebige $n_1,n_2\geq1$ mit $n_1\mid n_2$.
\end{theorem}
\begin{proof}
    Aus der Gruppentheorie wissen wir, dass endliche abelische Gruppen isomorph sind zu einer direkten Summe von zyklischen Gruppen \textbf{Quelle einfügen}, also
    \[\mathbb{Z}_{n_1}\oplus \mathbb{Z}_{n_2}\oplus \dots\oplus\mathbb{Z}_{n_r}\]
    mit $n_i\mid n_{i+1}$ für $i\geq 1$. Es gibt nun $n_1$ der Ordnung teilt $n_1$ Elemente für jedes i für die Gruppe $\mathbb{Z}_{n_i}$. Daher hat $E(\mathbb{F}_q)$ hat $n^r_1$ elemente der Ordnung teilt $n_1.$ Mit Theorem 3.2 sind nun $n^2_1$ solcher Punkte. Daher ist $r\leq 2$ und damit gezeigt.

\end{proof}



\subsection{Quadratische Reste und Legendre-Symbol, die Gruppen Ordnung bestimmensollte das nicht eher nach Hasse kommen?}

\cite{washington08}, S.104f.
Legendre Symbol nutzen, wenn man allgemien bestimmen will, ob es zu x einen quadratischen Rest mod q gibt, denn nur dann exestiert y und damit kann der Punkte auf der Kurve in einem endlichen Körper liegen.


Es ist also wichtig zu bestimmen ob es für ein $x$ auch ein $y$ gibt, also ob der berechnete Wert $a=y^2 \mod p$ ein Quadrat ist oder nicht. Hierfür können wir das Legendre-Symbol verwenden

\begin{definition}[Legendre-Symbol]
    Sei \(p\) eine ungerade Primzahl und \(a\in\mathbb Z\). Das Legendre-Symbol ist definiert durch

    \[
        \left(\frac{a}{p}\right)
        =
        \begin{cases}
            1,  & \text{falls } a \not\equiv 0 \pmod p \text{ ein quadratischer Rest modulo } p \text{ ist}, \\
            -1, & \text{falls } a \text{ ein quadratischer Nichtrest modulo } p \text{ ist},                 \\
            0,  & \text{falls } a \equiv 0 \pmod p.
        \end{cases}
    \]
\end{definition}

Es gibt zwei Methoden um das Legendre-Symbol effizient zu berechnen:
\begin{itemize}
    \item Das Eulersche Kriterium
    \item Das Quadratische Reziprozitätsgesetz
\end{itemize}
Auf beide wird nicht genauer eingeganngen, aber sie sind wichtig um Punkte für jede beliebige elliptische Kurve zu bestimmmen, da dadurch überprüft werden kann, ob es sich um einen quadratischen Rest oder nicht handelt.
Denn \(y^2=x^3+Ax+B\)
In \cite{muller} S.51ff. Kapitel 8 kann mehr zu dem Thema quadratische Reste nachgelesen werden und effiziente Methoden sowie die Beweise dazu gefunden werden.
Das Legendre-Symbol kann verallgemeinert werden über alle endlichen Köprer $\mathbb{F}_q$.
\begin{definition}
    (vgl. \cite{washington08},S. 104)
    Sei \(\mathbb F_q\) ein endlicher Körper mit \(q\) ungerade. Für \(x\in\mathbb F_q\) definieren wir

    \[
        \left(\frac{x}{\mathbb F_q}\right)
        =
        \begin{cases}
            1,  & \text{falls } t^2=x \text{ eine Lösung } t\in\mathbb F_q^\times \text{ besitzt}, \\
            -1, & \text{falls } t^2=x \text{ keine Lösung } t\in\mathbb F_q \text{ besitzt},       \\
            0,  & \text{falls } x=0.
        \end{cases}
    \]

\end{definition}
\begin{theorem}  (vgl. \cite{washington08},S.105)
    Sei $E$ eine elliptische Kurve definiert mit \(y^2=x^3+Ax+B\) über $\mathbb{F}_q$. Dann gilt
    \[
        |E(\mathbb{F}_q)|=q+1+\sum_{x\in\mathbb{F}_q}\left(\frac{x^3+Ax+B}{\mathbb{F}_q}\right)
    \]
\end{theorem}
\begin{proof}

\end{proof}

\section{Gruppenstruktur von \(E(\mathbb{F}_q)\)}

\begin{theorem}(vgl. \cite{silverman06}, S. 29)
    Über einem endlichen Körper ist die Gruppe der Punkte $E(\mathbb{F}_p)$ zusammen mit \(\infty\) immer eine zyklische Gruppe oder ein Produkt von zyklischen Untergruppen
\end{theorem}
Bevor wir das Theorem beweisen betrachten wir die Aussage an einem Beispiel:
\begin{beispiel}
    \cite{guneysu24}, S. 285f.
    Wir wollen alle Punkte auf der Kurve
    \[
        E:y^2\equiv x^3 +x+5 \pmod {11}
    \]
    Wir berechnen nun alle Punkte auf der Kurve, wir nehmen an das die Gruppe aller Punkte auf dieser Kurve zyklisch ist und außerdem ist die Ordnung $\left| E(\mathbb{F}_{11}) \right|=11$, Wir wählen den Punkt \(P=(0,4)\). Dieser liegt auf der Kurve, denn

    \[
        4^2 = 16 \equiv 5 \pmod{11}
    \]

    und

    \[
        0^3+0+5 = 5 \pmod{11}.
    \]

    Wir berechnen nun die Vielfachen von \(P\):

    \[
        \begin{aligned}
            2P  & = (0,4)+_E(0,4) = (5,5), \\
            3P  & = 2P+_EP = (10,5),       \\
            4P  & = (2, 9),                \\
            5P  & = (7, 6),                \\
            6P  & = (7, 5)                 \\
            7P  & = (2, 2),                \\
            8P  & = (10, 6),               \\
            9P  & = (5, 6),                \\
            10P & = (0, 7),                \\
            11P & = \infty,                \\
            12P & = (0,4),                 \\
            13P & = (5,5)
        \end{aligned}
    \]

    Wie man sieht beginnt nun die Aufzählung von neuem und es handelt sich daher um eine zyklische Gruppe.

    \begin{figure}[H]
        \centering
        \includegraphics[width=0.5\linewidth]{ell_curve_finite_field.png}
        \caption{Elliptische Kurve über $\mathbb{F}_5$, eigene Darstellung}
        \label{fig:placeholder}
    \end{figure}

    Die Punkte und das Bild wurden mit einem eigenen Python Skript generiert, siehe Anhang.
\end{beispiel}

Um das diskrete Logarithmus für ein Kryptosystem aufbzubauen, müssen wir die Ordnung von einer Gruppe kennen. Dies ist deutlich herausfordernder als das Berechnen des eben gezeigten Beispiels, bei dem extra eine passende Kurve gewählt wurde. Da zuerst überprüft werden muss ob zu einem x Wert der berechnete Wert auch einen quadratischen Rest hat und damit der Punkt exestiert.
\begin{theorem}
Sei $E$ eine elliptische Kurve über einem Körper $K$ und sei $n\in\mathbb N$.
Falls $\operatorname{char}(K)\nmid n$, dann gilt
\[
    E[n]\cong \mathbb Z/n\mathbb Z \oplus \mathbb Z/n\mathbb Z,
\]
wobei
\[
    E[n]=\{P\in E(\overline K)\mid [n]P=\mathcal O\}
\]
die Gruppe der $n$-Torsionspunkte bezeichnet.
\end{theorem}
vgl. \cite{washington08} Washington S.77+79

\begin{proof}
    Der Beweis hierzu ist sehr langwierig und findet sich daher in \cite{washington08} III Abschnitt 3.2.
\end{proof}


\section{Punktzählung}
Im vorangenganen Abschnitt wurde gezeigt wie man bei einfachen Gruppen vorgehen kann um die Ordnung von diesen zu bestimmen. Die Punktzählung für endliche Körper großer Primzahlordnung ist aber sehr unpraktikabel, denn es hat Komplexität \(\O(p^{1+\epsilon}, \epsilon>0)\)(vgl. \cite{schoof95},S. 2020). Daher soll in diesem Abschnitt Ideen entwickelt werden, wie man die Anzahl der Punkte einer Untergruppe von \(E(\mathbb{F}_q)\) berechnen kann.  Dies ist wichtig um später wichtig um passende Bedingungen für eine Kurve für ein Krpyotgraphisches Verfahren zu wählen. Hierfür wird vorallem der Algorithmus von Schoof und die Hasse Schranke benötigt. Der Beweis dieser benötigt eine Vielzahl von Eigenschaften die weit über die Arbeit hinausreichen, daher werden nur einige von ihnen aufgelistet, sodass zumindest eine Skizze des Beweises möglich ist. In \cite{silverman09} S. 138 findet man den ausführlichen Beweise. Dieser Abschnitt basier auf \cite{silverman09}, S.66f.




\subsection{Isogenien und Endmorphismen}
\begin{definition}[Isogenien]
    Seien $E_1$ und $E_2$ elliptische Kurven über
    einem Körper $K$. Eine Isogenie von $E_1$ nach $E_2$
    ist ein Morphismus
    \[
        \phi:E_1\to E_2
        \qquad\text{mit}\qquad
        \phi(\mathcal O_{E_1})=\mathcal O_{E_2}.
    \]
    Die Kurven $E_1$ und $E_2$ heißen isogen, wenn
    eine von null verschiedene Isogenie zwischen
    ihnen existiert, d.h. \(\phi(E_1)\neq\{\mathcal{O}\}\)
\end{definition}
\begin{beispiel}
    Für jedes $m\in\mathbb{Z}$ definiert man die Multiplikationsabbildung
    \[
        [m]: E\to E.
    \]
    Für $m>0$ gilt
    \[
        [m](P)=\underbrace{P+\dots+P}_{m\text{-mal}}.
    \]
    $[m]$ ein Morphismus ist. Für $m=1$ gilt
    $[1]=\operatorname{id}_E$, also ist $[1]$ ein Morphismus. Sei nun
    $[m]$ ein Morphismus. Dann gilt
    \[
        [m+1]=\mu\circ([m],\operatorname{id}_E),
    \]
    wobei $\mu:E\times E\to E$, $(P,Q)\mapsto P+Q$, nach
    \cite[III, Theorem 3.6, S.~64]{silverman09} ein Morphismus ist. Für $m<0$ gilt \[[m](P)=[-m](-P)\]. Da die Negationsabbildung ebenfalls ein Morphismus ist, ist $[m]$ auch für negative $m$ ein Morphismus.
    Schließlich gilt für jedes $m\in\mathbb Z$ die Gleichung
    $[m](\mathcal O)=\mathcal O$. Daher ist $[m]:E\to E$ eine Isogenie.
\end{beispiel}
Aus \cite{silverman09},S.67

\begin{definition}

    Seien \(E_1\) und \(E_2\) elliptische Kurven über einem Körper \(K\), und sei \(\phi:E_1\to E_2\) eine nichtkonstante rationale Abbildung, die über \(K\) definiert ist. Durch Komposition mit \(\phi\) wird eine Einbettung der zugehörigen Funktionenkörper induziert:
    \[
        \phi^*:K(E_2)\to K(E_1),\qquad
        f\mapsto f\circ\phi.
    \]
    Die Abbildung \(\phi\) heißt \textbf{separabel} beziehungsweise inseparabel, wenn die durch \(\phi\) induzierte Körpererweiterung
    \[
        K(E_1)/\phi^*(K(E_2))
    \]
    separabel beziehungsweise inseparabel ist.
    Für eine konstante Abbildung setzt man \(\deg(\phi)=0.\). Ist \(\phi\) nicht konstant, so definiert man den \textbf{Grad} durch
    \[
        \deg(\phi)
        =
        [K(E_1):\phi^*(K(E_2))].
    \]
    Der Grad von \(\phi\) entspricht somit dem Grad der endlichen Körpererweiterung \(K(E_1)/\phi^*(K(E_2))\).

\end{definition}

\begin{lemma}
    \label{lem:deg-isogenien}
    Sei $\phi:E_1\to E_2$ eine von null verschiedene seperable Isogenie zwischen zwei elliptischen Kurven über $K$. Dann gilt
    \[
        \#\ker \phi = \deg \phi,
    \]
    wobei der Kern auf den Punkten über $\overline{K}$ betrachtet wird.
\end{lemma}

\begin{proof}
    Der Beweis hierfür lässt sich in \cite{silverman09}, S.72f.  Theorem III.4.10. (c) finden.
\end{proof}

\subsection{Frobenius-Endomorphismus}
\begin{definition}[Frobenius Abbildung](vgl. \cite{washington08}, S. 98f.)
    \label{def:frobenius}
    Sei $E$ eine elliptische Kurve, gegeben durch eine
    Weierstraßgleichung über $\mathbb F_q$, wobei
    $q=p^r$ eine Primzahlpotenz ist. Die Abbildung
    \[
        \phi_q:E\to E,\qquad
        (x,y)\mapsto(x^q,y^q),\qquad
        \mathcal O\mapsto\mathcal O,
    \]
    heißt Frobenius-Endomorphismus.
\end{definition}

\begin{lemma}
    Sei $E$ eine elliptische Kurve über \(\mathbb{F}_q\). Dann ist $\phi_q$ ein Endomorphismus von $E$ mit dem Grad $q$ und $\phi_q$ ist nicht seperabel.
\end{lemma}
\begin{proof} (vlg. \cite{washington08}, S.52, eventuell etwas ausführlicher)
    Die Abbildung $\phi_q(x,y) = (x^q, y^q)$ ist wohldefiniert, da die Gleichung der elliptischen Kurve Koeffizienten in $\mathbb{F}_q$ besitzt und somit unter der Abbildung $x \mapsto x^q$ invariant ist.

    Da die Gruppenoperation auf $E$ durch rationale Funktionen mit Koeffizienten in $\mathbb{F}_q$ gegeben ist und $\phi_q$ ein Körperhomomorphismus ist, gilt
    \[
        \phi_q(P+Q) = \phi_q(P) + \phi_q(Q)
    \]
    für alle $P,Q \in E(\overline{\mathbb{F}}_q)$.
    Damit ist $\phi_q$ ein Endomorphismus von $E$.

    Der Grad von $\phi_q$ ist $q$.

    Da in Charakteristik $p$ gilt $q = p^r$ und die Ableitung von $x^q$ gleich null ist, ist $\phi_q$ nicht separabel.
\end{proof}

\begin{corollary}
    \label{korr:abbildung-seperabel}
    Sei $E$ eine elliptische Kurve über einem endlichen Körper $\mathbb{F}_q$ mit Charakterstik $p>0$. Dann gilt für den Frobenius-Endomorphismus $\phi:E\to E$ und $m,n\in\mathbb{Z}$, die Abbildung
    \[
        m+n\phi:E\to E
    \]
    ist genau dann seperabel, wenn $p \nmid m$. Daraus folgt die Abbildung $1-\phi$ ist seperabel.
\end{corollary}

\begin{proof}
    Den Beweis findet man in \cite{silverman09}, S.79.
\end{proof}

\begin{corollary}
    \label{kor:gradabbildung}
    Für die elliptische Kurven $E_1$ und $E_2$ ist die Abbildung
    \[
        \deg:\operatorname{Hom}(E_1,E_2)\to\mathbb{Z}, \qquad \phi\mapsto\deg(\phi),
    \]
    eine positivie definite quadratische Form.
\end{corollary}
\begin{proof}
    Ein Beweis findet sich in \cite{silverman09}, S. 85 Korrolar 6.3.
\end{proof}

\subsection{Hasse-Schranke}
\cite{werner02}, S.60--61 und \cite{guneysu24} S.287.
\begin{satz}[Satz von Hasse]
    \label{def:hasse-schranke}
    Sei $E$ eine elliptische Kurve über einem endlichen Kröper $\mathbb{F}_q$. Dann erfüllt die Ordnung von $E(\mathbb{F}_q)$,
    \[|\#E(\mathbb{F}_q)-q-1 |\leq 2\sqrt q\]
\end{satz}
Es ist nun möglich die Gruppenordnung von $E(\mathbb{F})$ rechteinfach abzuschätzen durch
\[
    -2\sqrt{q}+q+1\leq \#E(\mathbb{F})\leq 2\sqrt{q}+q+1.
\]
Dieses Theorem hat nun eine weitrechende Auswirkung auf die Implementierung von elliptischen Kurven, denn Hasses Schranke gibt an, dass wenn wir eine elliptische Kurven mit zum Beispiel $2^{256}$ Elementen benötigen. Benötigen wir eine Primzahl mit in etwa 256 Bit Länge vgl.\cite{guneysu24} S.287.

\begin{definition}[Bilinearform]
    \label{def:biliniearform}
    Sei $V$ ein Vektorraum über einem Körper $K$. Eine Abbildung
    \[
        \langle \cdot,\cdot\rangle:V\times V\to K
    \]
    heißt Bilinearform, falls sie in beiden Argumenten linear ist. Das heißt, für alle $x,y,z\in V$ und $\lambda,\mu\in K$ gilt
    \[
        \langle \lambda x+\mu y,z\rangle
        =
        \lambda\langle x,z\rangle+\mu\langle y,z\rangle
    \]
    sowie
    \[
        \langle x,\lambda y+\mu z\rangle
        =
        \lambda\langle x,y\rangle+\mu\langle x,z\rangle.
    \]
\end{definition}
\begin{satz}[Cauchy-Schwarzsche Ungleichung]
    Für eine Bilinearform auf einem Vektorraum $V$ und $x,y\in V$ gilt
    \[
        \langle x,y\rangle^2\leq \langle x,x\rangle\langle y,y\rangle
    \]
\end{satz}
\begin{proof}
    Ein Beweis findet sich in \cite{poeschel14} S.115.
\end{proof}

\begin{proof}[Hasse]
    Sei $E$ eine elliptische Kurve über \(\mathbb{F}_q\). Außerdem sei \(\phi: E\to E, \quad (x,y)\mapsto (x^q,y^q)\) ein Frobenius Endomorphismus. Für ein $P=(x,y)\in E(\overline{\mathbb{F}}_q$) gilt
    \[
        P\in E(\overline{\mathbb{F}}_q) \Leftrightarrow x,y\in \mathbb{F}_q.
    \]

    Das ist äquivalent zu
    \[
        x^q=x \quad \text{ und } \quad y^q=y.
    \]
    Also gilt
    \[
        \phi(P)=(x^q,y^q)=(x,y)=P.
    \]
    Folglich erhält man
    \[
        P\in E(\mathbb{F}_q), \Leftrightarrow \phi(P)=P.
    \]
    Betrachtet man nun den Endomorphismus
    \[1-\phi:E\to E.\] Für einen Punkt $P\in E(\overline{\mathbb{F}}_q)$ gilt
    \[
        (1-\phi)(P)=\mathcal O
        \Leftrightarrow
        P-\phi(P)=\mathcal O
        \Leftrightarrow
        P=\phi(P).
    \]
    Somit folgt
    \[
        E(\mathbb{F}_q)=\ker(1-\phi).
    \]
    Mit \ref{korr:abbildung-seperabel} folgt $1-\phi$ separabel ist, da p nicht 1 teilt. Nun kann man das Lemma \ref{lem:deg-isogenien} anwenden,
    \[
        \#E(\mathbb{F}_q)
        =\#\ker(1-\phi)
        =\deg(1-\phi).
    \]
    Nach Korollar \ref{kor:gradabbildung} ist die Gradabbildung \(\deg(1-\phi)\) eine positiv definite quadratische Form. Die zugehörige symmetrische Bilinearform ist definiert durch
    \[
        \langle \alpha,\beta\rangle
        =
        \deg(\alpha+\beta)-\deg(\alpha)-\deg(\beta).
    \]
    Es wird nun die Bilinearform verwendet, sodass gilt
    \begin{align*}
        \Leftrightarrow & \langle \alpha,\beta\rangle=\deg(\alpha+\beta)- \deg(\alpha)-\deg(\beta) \\
                        & \deg(\alpha+\beta)= \langle\alpha,\beta\rangle+\deg(\alpha)+\deg(\beta)
    \end{align*}
    Für den gegebenen Fall ist $\alpha=1$ und \(\beta =-\phi\), damit folgt
    \[
        \deg(1-\phi)=\langle 1,-\phi\rangle+\deg(1)+\deg(-\phi)+
    \]
    Es ist \(\deg(1)=1\) und \(\deg(-\phi)=\deg(\phi)=q\) und damit
    \[
        \deg(1-\phi)=\langle 1,-\phi\rangle+1+q
    \]
    mit Biliniearität gilt
    \[
        \deg(1-\phi)=1+q-\langle 1,\phi\rangle.
    \]
    Jetzt fehlt es noch zu zeigen, dass \(|\langle1,\phi\rangle|\leq 2\sqrt{q}\).
    Mit der Cauchy-Schwarz Ungleichung gilt \(|\langle\alpha,\beta\rangle|^2\leq \langle\alpha,\alpha\rangle\langle\beta,\beta\rangle\). Für $\alpha=1$ und \(\beta=\phi\) ergibt sich

    \[
        |
        \langle1,\phi\rangle|^2\leq \langle1,1\rangle\langle\phi,\phi\rangle.
    \]
    Es gilt nach Definition
    \[
        \langle 1,1\rangle = \deg(1+1)-\deg(1) - \deg(1)= 2
    \]
    und
    \[
        \langle\phi,\phi\rangle=\deg(\phi+\phi)-2\deg(\phi)
    \]
    mit \[
        \deg(\phi+\phi)
        =\deg([2]\circ\phi)
        =\deg([2])\deg(\phi)
        =4q
    \]
    folgt
    \[
        \langle\phi,\phi\rangle
        =\deg(\phi+\phi)-2\deg(\phi)
        =4q-2q
        =2q.
    \]
    Nun setzt man noch alles in Cauchy-Schwarz ein, dann folgt
    \[
        |\langle 1,\phi\rangle|^2\leq 2\cdot 2q = 4q
    \]
    und zieht man nun die Wurzel gilt
    \[
        |\langle 1,\phi\rangle|\leq 2\sqrt{q}
    \]
    und damit gilt die gewünschte Aussage.
\end{proof}
\cite{silverman09}, S.138
\cite{silverman06}, S.100






\section{Double-and-Add-Algorithmus}
\label{sec:double-and-add}

Dieser Abschnitt basiert auf \cite{silverman06}, S.30--31 und \cite{buell24}, S.217--218.

Sei $P$ ein Punkt einer additiv geschriebenen Gruppe und
$n \in \mathbb{N}$. Ziel ist es, das skalare Vielfache
\[
    nP = \underbrace{P + P + \dots + P}_{n\text{-mal}}
\]
effizient zu berechnen.

Eine naive Berechnung erfolgt durch sukzessive Addition:
\begin{align*}
    2P & = P + P,      \\
    3P & = 2P + P,     \\
       & \ \vdots      \\
    nP & = (n-1)P + P.
\end{align*}
Hierfür werden insgesamt $n-1$ Gruppenadditionen benötigt. In
Abhängigkeit vom Zahlenwert $n$ liegt die Laufzeit dieser Methode daher in $\mathcal{O}(n)$. Für die Betrachtung der tatsächlichen
Eingabegröße ist jedoch die Anzahl der Bits entscheidend, die zur Darstellung von $n$ benötigt wird. Bezeichnet \(m = \lfloor \log_2(n) \rfloor + 1\) die Bitlänge von $n$, so gilt \( 2^{m-1} \leq n < 2^m \). Der Zahlenwert $n$ wächst somit exponentiell mit seiner Bitlänge
$m$. Daher kann die Laufzeit der naiven Methode äquivalent als \( \mathcal{O}(2^m)\) angegeben werden. Insbesondere bei kryptographischen Verfahren wie dem Diffie-Hellman-Schlüsselaustausch aus
Kapitel 5 ~\ref{sec:diffie-hellman} müssen skalare Vielfache für sehr große Skalare berechnet werden. Die naive Methode ist dafür
nicht praktikabel. Eine wesentlich effizientere Berechnung
ermöglicht der Double-and-Add-Algorithmus.
\cite{guneysu24}, S. 288
\begin{definition}[Double-and-Add-Algorithmus]
    Sei $P$ ein Punkt auf einer elliptischen Kurve $E$ und $n\in \mathbb{N}$. Es soll nun das Vielfache $Q=np$ berechnet werden. Sei \( n = (b_{m-1}b_{m-2}\dots b_0)_2\qquad\text{mit}\qquad b_{m-1}=1\)
    die Binärdarstellung von $n$.Der Double-and-Add-Algorithmus
    verarbeitet die Bits von links nach rechts, beginnend mit dem
    höchstwertigen Bit $b_{m-1}$.
    \begin{enumerate}
        \item Initialisiere $Q = \mathcal{O}$
        \item Führe für $i=m-1,m-2,\dots,0$ die folgenden Schritte aus:
              \begin{enumerate}
                  \item Verdopple den aktuellen Punkt: \( Q = 2Q\)
                  \item Falls $b_i=1$, dann addiere $P$: \(Q=Q+P\)
              \end{enumerate}
    \end{enumerate}
    Nach der Verarbeitung aller Bits gilt $Q=nP$.
\end{definition}
Um die Komplexität zu bestimmen, schreibt man das Skalar $n$ nochmal mit Hilfe der Binärdarstellung
Aus der Binärdarstellung $n=n_0+n_1\cdot 2\cdot 2+n_2\cdot 2^2+\cdots +n_r\cdot 2^r$ mit $n_0,\dots,n_r\in\{0,1\}.$. Daraus folgt für das skalare Vielfache \(nP=n_0P+n_1\cdot2P\cdot 2+n_2\cdot 2^2P+\cdots +N_r\cdot 2^rP\).Der Punkt $2^kP$ wird durch $k$ aufeinanderfolgende Punktverdopplungen berechnet. Da $r=\lfloor\log_2(m)\rfloor$ gilt, benötigt der Double-and-Add-Algorithmus genau $\lfloor\log_2(m)\rfloor$ Verdopplungen. Bei annähernd gleichverteilten Bits kommen im Mittel etwa \(\frac{1}{2}\lfloor\log_2(m)\rfloor\) Punktadditionen hinzu. Die Laufzeit liegt somit in
$\mathcal{O}(\log m)$.
\begin{beispiel}
    Es soll das skalare Vielfache $46P$ berechnet werden. Zunächst
    wird der Skalar in Binärdarstellung überführt \( 46=(101110)_2\).
    Die Verarbeitung beginnt beim höchstwertigen Bit (\emph{most significant bit}, MSB). Da dieses stets gleich $1$ ist, wird der Startwert $Q=P$ gesetzt. Anschließend werden die
    verbleibenden Bits von links nach rechts verarbeitet.
    \begin{table}[htbp]
        \centering
        \begin{tabular}{c l l}
            \hline
            \textbf{Bit} & \textbf{Operation}       & \textbf{Zwischenergebnis} \\
            \hline
            $1$          & Initialisierung          & $Q=P$                     \\
            $0$          & Verdopplung              & $Q=2P$                    \\
            $1$          & Verdopplung und Addition & $Q=5P$                    \\
            $1$          & Verdopplung und Addition & $Q=11P$                   \\
            $1$          & Verdopplung und Addition & $Q=23P$                   \\
            $0$          & Verdopplung              & $Q=46P$                   \\
            \hline
        \end{tabular}
        \caption{Berechnung von $46P$ mit dem Double-and-Add-Algorithmus}
        \label{tab:double-and-add-26}
    \end{table}

    Für die Berechnung von $46P$ werden somit fünf Verdopplungen und
    drei zusätzliche Additionen benötigt. Insgesamt sind dies sechs
    Gruppenoperationen. Die naive Methode würde dagegen
    $46-1=45$ Additionen erfordern.
\end{beispiel}



\subsection{Baby steps and giant steps}
\cite{schoof95}, S.220 Komplexität ist \(\O(p^{1/4+\epsilon}), \epsilon>0\). Praktisch für großere Primzahlen, aber unpraktisch für Zahlen größer 20 Dezimalstellen

Es wird Hasse benötigt.
\subsection{Schofs Algorithm}

